\documentclass[11pt,a4paper]{amsart}
\usepackage[a4paper,margin=30mm]{geometry}
\usepackage[T1]{fontenc}
\usepackage{lmodern,amssymb,mathtools,microtype}
\usepackage[hidelinks]{hyperref}
\hypersetup{pdftitle={Minimizing the variance of interpoint distances in the plane},pdfauthor={Lazar Korsic}}
\newtheorem{theorem}{Theorem}[section]
\newtheorem{lemma}[theorem]{Lemma}
\newtheorem{corollary}[theorem]{Corollary}
\theoremstyle{remark}
\newtheorem{remark}[theorem]{Remark}
\newcommand{\E}{\mathbb E}
\newcommand{\R}{\mathbb R}
\newcommand{\Var}{\operatorname{Var}}
\newcommand{\ind}{\mathbf{1}}
\title{Minimizing the variance of interpoint distances in the plane}
\author{Lazar Korsi\'c}
\date{November 25, 2026}
\subjclass[2020]{60D05, 52A40, 26D15}
\keywords{Geometric probability, interpoint distance, variance, bounded density, negative type}
\begin{document}
\begin{abstract}
We prove that, among measurable planar sets of prescribed positive
finite area, the disk uniquely minimizes the variance of the distance
between two independent uniform points. More generally, among all
planar probability densities bounded above by a given constant, the
uniform density on a disk uniquely minimizes this variance, up to
translation and almost-everywhere equality. The proof combines the
negative type of Euclidean distance, a comparison with the potential
of a uniform disk, and an elementary universal mean-distance bound.
\end{abstract}
\maketitle

\section{Introduction}
In August 2021 a question on Math Stack Exchange \cite{MSE} asked
which planar region of prescribed area minimizes the variance of the
Euclidean distance between two independent uniform points, and
suggested that a disk is optimal. We establish this assertion without
any convexity, connectedness, or boundary regularity assumption. In
fact, the same conclusion holds in a larger class of probability
measures.

\begin{theorem}\label{th:main}
Let $M>0$, and let $X,Y$ be independent identically distributed random
vectors in $\R^2$ whose common probability density $f$ satisfies
$0\le f\le M$ almost everywhere. Then
\begin{equation}\label{eq:main}
 \Var(|X-Y|)\ge \frac{1}{M}
 \left(\frac1\pi-\frac{16384}{2025\pi^3}\right).
\end{equation}
Equality holds if and only if $f=M\ind_B$ almost everywhere for a disk
$B$ of area $1/M$.
\end{theorem}

\begin{corollary}\label{cor:set}
Let $A\subset\R^2$ be Lebesgue measurable with $0<|A|<\infty$,
and let $X,Y$ be independent uniform points in $A$. Then
\[
 \Var(|X-Y|)\ge |A|\left(\frac1\pi-
 \frac{16384}{2025\pi^3}\right).
\]
Equality holds if and only if $A$ agrees almost everywhere with a disk.
\end{corollary}
\begin{proof}
Apply Theorem~\ref{th:main} to $f=\ind_A/|A|$.
\end{proof}

The related problem of maximizing the mean distance under a constraint
on the second moment, but with no density bound, was solved by Buja,
Logan, Reeds and Shepp \cite{BLRS}. Its planar optimizer is the
projection of a uniform distribution on a sphere in $\R^3$ and has
an unbounded planar density. This differs from the constrained
minimization considered here. Mean-distance inequalities for convex
bodies have also been investigated by Bonnet, Gusakova, Th\"ale and
Zaporozhets \cite{BGTZ}. The present proof is self-contained; in
particular, it does not require the sharp bound of \cite{BLRS}.

\paragraph{Provenance and verification status.}
I asked OpenAI's ChatGPT models to investigate the problem, and suggested
initial experiments with annuli and a possible approach based on unions
of small disks. The models developed the proof, its bounded-density
extension, and the manuscript. This note belongs to the AI-generated
research collection. The argument has been rechecked by AI but has not
been independently refereed or formally verified. This draft precedes
completion of my own review. No claim of priority over all existing
literature is made.

\section{Preliminaries}
By translation we may assume that the common law $P$ of $X,Y$ has
centroid zero, whenever the second moment is finite. Write
\[
 I(P)=\int |x|^2\,dP(x),\qquad
 m(P)=\iint |x-y|\,dP(x)\,dP(y).
\]
Independence and the zero centroid give
\begin{equation}\label{eq:variance}
 \Var(|X-Y|)=2I(P)-m(P)^2.
\end{equation}
It suffices to consider finite second moments. Indeed, if
$\E|X|^2=\infty$, then for every fixed $y$ the inequality
$|x|^2\le 2|x-y|^2+2|y|^2$ implies
$\int |x-y|^2\,dP(x)=\infty$. Hence $\E|X-Y|^2=\infty$ by Tonelli's
 theorem. Such laws have no finite distance variance; throughout,
we may take the extended variance of a nonnegative random variable $D$
to mean $\inf_{c\in\R}\E(D-c)^2$, which is infinite in this case.

\begin{lemma}[Negative type]\label{lem:negative}
For probability measures $P,Q$ on $\R^2$ with finite first moments,
let $U_Q(x)=\int |x-y|\,dQ(y)$. Then
\begin{equation}\label{eq:negative}
 m(P)-m(Q)\le 2\int U_Q\,d(P-Q).
\end{equation}
\end{lemma}
\begin{proof}
For a finite signed measure $\sigma$ on $\R$ of total mass zero and
finite first absolute moment, the identity
\[
 |s-t|=\int_{\R}(\ind_{\{s\le z\}}-\ind_{\{t\le z\}})^2\,dz
\]
and Fubini's theorem yield
\[
 \iint |s-t|\,d\sigma(s)d\sigma(t)
 =-2\int_{\R}\sigma(({-\infty},z])^2\,dz\le0.
\]
The absolute integrability required for Fubini follows by integrating
$|s-t|$ against $|\sigma|\otimes|\sigma|$.
For $v\in\R^2$, with $e_\theta=(\cos\theta,\sin\theta)$,
\[
 |v|=\frac12\int_0^\pi |v\cdot e_\theta|\,d\theta.
\]
Applying the one-dimensional identity to the pushforward of $P-Q$
under each projection and integrating in $\theta$ proves
\[
 \iint |x-y|\,d(P-Q)(x)d(P-Q)(y)\le0.
\]
Expanding the left-hand side gives \eqref{eq:negative}.
\end{proof}

\section{The uniform disk and its potential}
Let $Q_t$ denote uniform probability measure on the centered disk
$B_t$ of radius $t>0$. Put
\[
 a=\frac{128}{45\pi}.
\]
The disk moments are
\begin{equation}\label{eq:diskmom}
 I(Q_t)=\frac{t^2}{2},\qquad m(Q_t)=at.
\end{equation}
For completeness, the first identity is immediate by polar integration.
For the second it suffices by scaling to take $t=1$. The overlap area
of two unit disks with centers distance $r\in[0,2]$ apart is
\[
 C(r)=2\arccos(r/2)-\frac r2\sqrt{4-r^2}.
\]
The difference of two independent uniform unit-disk points has density
$C(|z|)/\pi^2$ on $B_2$. Since $C'(r)=-\sqrt{4-r^2}$ for $0<r<2$,
integration by parts gives
\[
 m(Q_1)=\frac2\pi\int_0^2 r^2C(r)\,dr
 =\frac{2}{3\pi}\int_0^2 r^3\sqrt{4-r^2}\,dr
 =\frac{128}{45\pi}.
\]

\begin{lemma}[Disk potential]\label{lem:potential}
Write $U_{Q_t}(x)=u_t(|x|)$. Then, for all $r\ge0$,
\begin{equation}\label{eq:radial}
 u_t'(r)\le \frac rt,\qquad
 u_t(r)\le \frac{2t}{3}+\frac{r^2}{2t}.
\end{equation}
In particular, $r\mapsto u_t(r)-r^2/(2t)$ is nonincreasing.
\end{lemma}
\begin{proof}
The distributional identity $\Delta |x|=1/|x|$ holds on $\R^2$:
away from zero it follows by differentiation, and the boundary term
on a circle of radius $\varepsilon$ tends to zero as
$\varepsilon\downarrow0$. Convolving with the disk density gives
\[
 \Delta U_{Q_t}(x)=\frac{1}{\pi t^2}\int_{B_t}\frac{dy}{|x-y|}
 \le \frac{1}{\pi t^2}\int_{B_t}\frac{dy}{|y|}=\frac2t.
\]
For the inequality, translate the domain of integration by $-x$ and
compare with $B_t$. The part outside $B_t$ has the same area as the
missing part inside $B_t$; the integrand is at most $1/t$ on the
former and at least $1/t$ on the latter.
The second distributional derivatives of $|x|$ are locally integrable
and have magnitude bounded by a constant times $1/|x|$. Their
convolutions with the bounded, compactly supported disk density are
continuous, by continuity of translations in local $L^1$. Thus the
potential is $C^2$, and radial differentiation is justified. In
particular $u_t'(0)=0$. For $r>0$ we have
\[
 (ru_t'(r))'=r\Delta U_{Q_t}(x)\le 2r/t,
 \qquad |x|=r.
\]
Integration from zero proves the derivative bound. The other bound
follows by integration once more, using $u_t(0)=2t/3$.
\end{proof}

\section{Two comparisons}
Normalize $M=1$ for now. Put $R=\pi^{-1/2}$, $Q=Q_R$, and write
$I=I(P)$, $m=m(P)$, $I_0=I(Q)$ and $m_0=m(Q)$, with $P$ centered.
Thus $Q$ has density $\ind_{B_R}$, $I_0=R^2/2$ and $m_0=aR$.

\begin{lemma}[Comparison with the disk]\label{lem:diskcompare}
If $P$ has density $0\le f\le1$, then
\begin{equation}\label{eq:inertia}
 I\ge I_0,
\end{equation}
with equality only for $f=\ind_{B_R}$ almost everywhere. Moreover,
\begin{equation}\label{eq:comparison}
 I-I_0\ge R(m-m_0).
\end{equation}
\end{lemma}
\begin{proof}
Let $g=f-\ind_{B_R}$. It has integral zero, is nonpositive inside
$B_R$, and is nonnegative outside. Consequently
\[
 I-I_0=\int (|x|^2-R^2)g(x)\,dx\ge0.
\]
The integrand is nonnegative everywhere and strictly positive wherever
$g\ne0$ off the boundary circle, proving the equality assertion.
By Lemma~\ref{lem:potential}, the function
$\phi(r)=u_R(r)-r^2/(2R)$ is nonincreasing. Hence
\[
 \int\phi(|x|)g(x)\,dx
 =\int(\phi(|x|)-\phi(R))g(x)\,dx\le0.
\]
Lemma~\ref{lem:negative} now gives
\[
 m-m_0\le 2\int u_R(|x|)g(x)\,dx
 \le \frac{I-I_0}{R},
\]
as required.
\end{proof}

\begin{lemma}[Universal bound]\label{lem:universal}
Let $P$ be any centered probability measure on $\R^2$ with finite
second moment. Set $b=4/3-a$. Then
\begin{equation}\label{eq:universal}
 m(P)^2\le 4bI(P).
\end{equation}
In particular, since $0<b<1/2$,
\begin{equation}\label{eq:universalV}
 \Var(|X-Y|)\ge \left(\frac1{2b}-1\right)m(P)^2.
\end{equation}
\end{lemma}
\begin{proof}
Compare $P$ with $Q_t$ for arbitrary $t>0$. By
Lemmas~\ref{lem:negative} and \ref{lem:potential} and
\eqref{eq:diskmom},
\begin{align*}
 m(P)&\le 2\int U_{Q_t}(x)\,dP(x)-m(Q_t)\\
 &\le \frac{I(P)}t+\left(\frac43-a\right)t.
\end{align*}
If $I(P)>0$, take $t=\sqrt{I(P)/b}$. If $I(P)=0$, the law is a
point mass and the bound holds trivially. The stated range of $b$
follows, for example, from $3<\pi<22/7$. Finally apply
\eqref{eq:variance}.
\end{proof}

\begin{remark}
The universal bound is not sharp. The result of \cite{BLRS} implies
the stronger planar inequality $m(P)^2\le (\pi^2/6)I(P)$. The weaker
bound \eqref{eq:universal} already suffices for the argument below.
\end{remark}

\section{Proof of the main theorem}
For the normalized problem $M=1$, the disk has variance
\begin{equation}\label{eq:V0}
 V_0=2I_0-m_0^2=R^2(1-a^2)
 =\frac1\pi-\frac{16384}{2025\pi^3}.
\end{equation}
Write $V=2I-m^2$. We distinguish three ranges of $m$.

\smallskip\noindent
\emph{Case 1: $m\le m_0$.}
By \eqref{eq:inertia},
\[
 V\ge 2I_0-m^2\ge 2I_0-m_0^2=V_0.
\]
Equality forces $I=I_0$, hence $f=\ind_{B_R}$ almost everywhere.

\smallskip\noindent
\emph{Case 2: $m_0<m<2R-m_0$.}
By \eqref{eq:comparison},
\begin{align*}
 V-V_0&=2(I-I_0)-(m^2-m_0^2)\\
 &\ge (m-m_0)(2R-m_0-m)>0.
\end{align*}

\smallskip\noindent
\emph{Case 3: $m\ge 2R-m_0$.}
The universal bound \eqref{eq:universalV} gives
\[
 V\ge R^2\left(\frac1{2b}-1\right)(2-a)^2
 >R^2(1-a^2)=V_0.
\]
To verify the strict comparison, multiply the difference of the
coefficients of $R^2$ by $2b>0$ and substitute $b=4/3-a$:
\[
 (1-2b)(2-a)^2-2b(1-a^2)
 =\frac{-21a^2+50a-28}{3}.
\]
The bounds $3<\pi<22/7$ imply $448/495<a<1$. The polynomial
$p(a)=-21a^2+50a-28$ is increasing on this interval and satisfies
$p(448/495)=4172/81675>0$.

\smallskip
Since $a<1$, the three ranges cover every possible $m$. This proves
the theorem for $M=1$, including the equality assertion. For general
$M>0$, put $X'=\sqrt M X$, $Y'=\sqrt M Y$. The density of $X'$ is
$M^{-1}f(x/\sqrt M)\le1$, and
$\Var(|X'-Y'|)=M\Var(|X-Y|)$. Applying the normalized result and
scaling the equality case back proves Theorem~\ref{th:main}.

\begin{thebibliography}{9}
\bibitem{MSE}
Shaq (Math Stack Exchange username),
\emph{Minimize the variance of distances between 2 points for planar
region of given area}, Math Stack Exchange, August 14, 2021.
\url{https://math.stackexchange.com/questions/4224464/}.

\bibitem{BLRS}
A.~Buja, B.~F.~Logan, J.~A.~Reeds and L.~A.~Shepp,
\emph{Inequalities and positive-definite functions arising from a
problem in multidimensional scaling},
Ann. Statist. \textbf{22} (1994), no.~1, 406--438.
\href{https://doi.org/10.1214/aos/1176325376}{doi:10.1214/aos/1176325376}.

\bibitem{BGTZ}
G.~Bonnet, A.~Gusakova, C.~Th\"ale and D.~Zaporozhets,
\emph{Sharp inequalities for the mean distance of random points in
convex bodies}, Adv. Math. \textbf{386} (2021), 107813.
\href{https://doi.org/10.1016/j.aim.2021.107813}{doi:10.1016/j.aim.2021.107813}.
\end{thebibliography}
\end{document}
